% GENERATED FILE. Edit canonical lessons, then run npm run generate:books.
% canonical-source-hash: fnv1a64:8089ff88dd0df811
% canonical-lessons: TE-C368-sasanam, TE-C368-talapatram, TE-C368-caritrakarudu, TE-C368-nagarikata, TE-C368-yugam

\chapter{Written in Stone}
\label{ch:te-written-in-stone}
\hlchaptermodality{368}

\begin{chapteropening}
\textbf{By the end of this chapter:} \emph{I can say an inscription, a palm-leaf manuscript, a historian, a civilisation and an era.}

Complete the last lesson of chapter 368: I can say an inscription, a palm-leaf manuscript, a historian, a civilisation and an era.
\end{chapteropening}

\section[\texorpdfstring{\te{శాసనం} (śāsanaṁ)}{శాసనం (śāsanaṁ)}]{\te{శాసనం} (śāsanaṁ) — an inscription; an edict}

\label{lesson:TE-C368-sasanam}



\begin{quote}
Before the new one: say the Telugu for remains, then the Telugu for to unearth.
\end{quote}

\subsection*{You'll want to know: \te{శాసనం}}
\textbf{\te{శాసనం}} — \emph{śāsanaṁ} — \textquotedblleft{}an inscription; an edict\textquotedblright{}.

\begin{quote}
\te{ఆ} \te{రాతి} \te{మీద} \te{ఒక} \te{పాత} \te{శాసనం} \te{ఉంది}.
\end{quote}

\begin{quote}
\emph{ā rāti mīda oka pāta śāsanaṁ undi.}
\end{quote}

\begin{quote}
There is an old inscription on that stone.
\end{quote}

Kings carved their orders and gifts on stone and copper.

\begin{cousinweb}
\textbf{\te{శాసనం}} (\emph{śāsanaṁ}) is borrowed from Sanskrit \emph{śāsana}, \textquotedblleft{}a command; a royal edict\textquotedblright{} (CDIAL 12420).
\end{cousinweb}

\begin{grammarlens}[title={What you've built}]
One more word for telling what happened.
\end{grammarlens}

\subsection*{Your turn}
\begin{itemize}
\raggedright
  \item \emph{Say it:} \emph{śāsanaṁ}
  \item \emph{Say it:} \emph{śāsanaṁ}, once more
  \item \emph{Say it:} \emph{ā rāti mīda oka pāta śāsanaṁ undi.} — There is an old inscription on that stone.
  \item \emph{Recall:} say the Telugu for remains, then the Telugu for to unearth, then say \emph{śāsanaṁ} again
\end{itemize}

\begin{tcolorbox}[breakable,colback=teal!4,colframe=teal!35!black,title={Before you move on}]
What does \te{శాసనం} mean? (\textquotedblleft{}An inscription; an edict\textquotedblright{}.) Say it once more.
\end{tcolorbox}
\section[\texorpdfstring{\te{తాళపత్రం} (tāḷapatraṁ)}{తాళపత్రం (tāḷapatraṁ)}]{\te{తాళపత్రం} (tāḷapatraṁ) — a palm-leaf manuscript}

\label{lesson:TE-C368-talapatram}



\begin{quote}
Before the new one: say the Telugu for to unearth, then the Telugu for an inscription.
\end{quote}

\subsection*{You'll want to know: \te{తాళపత్రం}}
\textbf{\te{తాళపత్రం}} — \emph{tāḷapatraṁ} — \textquotedblleft{}a palm-leaf manuscript\textquotedblright{}.

\begin{quote}
\te{మా} \te{తాతగారి} \te{దగ్గర} \te{పాత} \te{తాళపత్రాలు} \te{ఉన్నాయి}.
\end{quote}

\begin{quote}
\emph{mā tātagāri daggara pāta tāḷapatrālu unnāyi.}
\end{quote}

\begin{quote}
Grandfather has old palm-leaf manuscripts.
\end{quote}

Plural: \te{తాళపత్రాలు}, \emph{tāḷapatrālu}.

\te{తాళ}, the palmyra palm, and \te{పత్రం}, a leaf: the books of the days before paper.

\begin{cousinweb}
\textbf{\te{తాళపత్రం}} (\emph{tāḷapatraṁ}) is borrowed from Sanskrit: \emph{tāla}, \textquotedblleft{}the palmyra palm\textquotedblright{}, and \emph{pattra}, \textquotedblleft{}a leaf\textquotedblright{} (Monier-Williams).
\end{cousinweb}

\begin{grammarlens}[title={What you've built}]
One more word for telling what happened.
\end{grammarlens}

\subsection*{Your turn}
\begin{itemize}
\raggedright
  \item \emph{Say it:} \emph{tāḷapatraṁ}
  \item \emph{Say it:} \emph{tāḷapatraṁ}, once more
  \item \emph{Say it:} \emph{mā tātagāri daggara pāta tāḷapatrālu unnāyi.} — Grandfather has old palm-leaf manuscripts.
  \item \emph{Recall:} say the Telugu for to unearth, then the Telugu for an inscription, then say \emph{tāḷapatraṁ} again
\end{itemize}

\begin{tcolorbox}[breakable,colback=teal!4,colframe=teal!35!black,title={Before you move on}]
What does \te{తాళపత్రం} mean? (\textquotedblleft{}A palm-leaf manuscript\textquotedblright{}.) Say it once more.
\end{tcolorbox}
\section[\texorpdfstring{\te{చరిత్రకారుడు} (caritrakāruḍu)}{చరిత్రకారుడు (caritrakāruḍu)}]{\te{చరిత్రకారుడు} (caritrakāruḍu) — a historian}

\label{lesson:TE-C368-caritrakarudu}



\begin{quote}
Before the new one: say the Telugu for an inscription, then the Telugu for a palm-leaf manuscript.
\end{quote}

\subsection*{You'll want to know: \te{చరిత్రకారుడు}}
\textbf{\te{చరిత్రకారుడు}} — \emph{caritrakāruḍu} — \textquotedblleft{}a historian\textquotedblright{}.

\begin{quote}
\te{చరిత్రకారుడు} \te{ఆ} \te{శాసనం} \te{చదివాడు}.
\end{quote}

\begin{quote}
\emph{caritrakāruḍu ā śāsanaṁ cadivāḍu.}
\end{quote}

\begin{quote}
The historian read the inscription.
\end{quote}

Plural: \te{చరిత్రకారులు}, \emph{caritrakārulu}.

\te{చరిత్ర}, history, with -\te{కారుడు}, one who makes: compare \te{కళాకారుడు}, an artist.

\begin{cousinweb}
\textbf{\te{చరిత్రకారుడు}} (\emph{caritrakāruḍu}) is built on Sanskrit \emph{caritra}, \textquotedblleft{}going; behaviour\textquotedblright{} (CDIAL 4691), the source of \te{చరిత్ర}, with \emph{-kāruḍu}, \textquotedblleft{}one who does\textquotedblright{}.
\end{cousinweb}

\begin{grammarlens}[title={What you've built}]
One more word for telling what happened.
\end{grammarlens}

\subsection*{Your turn}
\begin{itemize}
\raggedright
  \item \emph{Say it:} \emph{caritrakāruḍu}
  \item \emph{Say it:} \emph{caritrakāruḍu}, once more
  \item \emph{Say it:} \emph{caritrakāruḍu ā śāsanaṁ cadivāḍu.} — The historian read the inscription.
  \item \emph{Recall:} say the Telugu for an inscription, then the Telugu for a palm-leaf manuscript, then say \emph{caritrakāruḍu} again
\end{itemize}

\begin{tcolorbox}[breakable,colback=teal!4,colframe=teal!35!black,title={Before you move on}]
What does \te{చరిత్రకారుడు} mean? (\textquotedblleft{}A historian\textquotedblright{}.) Say it once more.
\end{tcolorbox}
\section[\texorpdfstring{\te{నాగరికత} (nāgarikata)}{నాగరికత (nāgarikata)}]{\te{నాగరికత} (nāgarikata) — a civilisation}

\label{lesson:TE-C368-nagarikata}



\begin{quote}
Before the new one: say the Telugu for a palm-leaf manuscript, then the Telugu for a historian.
\end{quote}

\subsection*{You'll want to know: \te{నాగరికత}}
\textbf{\te{నాగరికత}} — \emph{nāgarikata} — \textquotedblleft{}a civilisation\textquotedblright{}.

\begin{quote}
\te{ఆ} \te{నది} \te{ఒడ్డున} \te{ఒక} \te{పురాతన} \te{నాగరికత} \te{ఉండేది}.
\end{quote}

\begin{quote}
\emph{ā nadi oḍḍuna oka purātana nāgarikata uṇḍēdi.}
\end{quote}

\begin{quote}
An ancient civilisation once stood on the bank of that river.
\end{quote}

Built on \te{నగరం}, a city: the life of cities.

\begin{cousinweb}
\textbf{\te{నాగరికత}} (\emph{nāgarikata}) is built on Sanskrit \emph{nāgarika}, \textquotedblleft{}born or living in a town, civic\textquotedblright{} (Monier-Williams), from \emph{nagara}, \textquotedblleft{}a city\textquotedblright{}, with the noun ending \emph{-ta}.
\end{cousinweb}

\begin{grammarlens}[title={What you've built}]
One more word for telling what happened.
\end{grammarlens}

\subsection*{Your turn}
\begin{itemize}
\raggedright
  \item \emph{Say it:} \emph{nāgarikata}
  \item \emph{Say it:} \emph{nāgarikata}, once more
  \item \emph{Say it:} \emph{ā nadi oḍḍuna oka purātana nāgarikata uṇḍēdi.} — An ancient civilisation once stood on the bank of that river.
  \item \emph{Recall:} say the Telugu for a palm-leaf manuscript, then the Telugu for a historian, then say \emph{nāgarikata} again
\end{itemize}

\begin{tcolorbox}[breakable,colback=teal!4,colframe=teal!35!black,title={Before you move on}]
What does \te{నాగరికత} mean? (\textquotedblleft{}A civilisation\textquotedblright{}.) Say it once more.
\end{tcolorbox}
\section[\texorpdfstring{\te{యుగం} (yugaṁ)}{యుగం (yugaṁ)}]{\te{యుగం} (yugaṁ) — an era; an age}

\label{lesson:TE-C368-yugam}



\begin{quote}
Before the new one: say the Telugu for a historian, then the Telugu for a civilisation.
\end{quote}

\subsection*{You'll want to know: \te{యుగం}}
\textbf{\te{యుగం}} — \emph{yugaṁ} — \textquotedblleft{}an era; an age\textquotedblright{}.

\begin{quote}
\te{ఇది} \te{కంప్యూటర్ల} \te{యుగం}.
\end{quote}

\begin{quote}
\emph{idi kampyūṭarla yugaṁ.}
\end{quote}

\begin{quote}
This is the age of computers.
\end{quote}

Plural: \te{యుగాలు}, \emph{yugālu}.

The old stories count time in \te{యుగాలు}, ages; today the word also names an era of history.

\begin{cousinweb}
\textbf{\te{యుగం}} (\emph{yugaṁ}) is borrowed from Sanskrit \emph{yuga} (CDIAL 10482), whose first sense is \textquotedblleft{}a yoke\textquotedblright{}; it came to mean an age of the world.
\end{cousinweb}

\begin{grammarlens}[title={What you've built}]
One more word for telling what happened.
\end{grammarlens}

\subsection*{Your turn}
\begin{itemize}
\raggedright
  \item \emph{Say it:} \emph{yugaṁ}
  \item \emph{Say it:} \emph{yugaṁ}, once more
  \item \emph{Say it:} \emph{idi kampyūṭarla yugaṁ.} — This is the age of computers.
  \item \emph{Recall:} say the Telugu for a historian, then the Telugu for a civilisation, then say \emph{yugaṁ} again
\end{itemize}

\begin{tcolorbox}[breakable,colback=teal!4,colframe=teal!35!black,title={Before you move on}]
What does \te{యుగం} mean? (\textquotedblleft{}An era; an age\textquotedblright{}.) Say it once more.
\end{tcolorbox}
